% Generated draft: requires mathematical and source review.
\documentclass[11pt]{article}
\usepackage[margin=25mm]{geometry}
\usepackage{fontspec}
\setmainfont{Latin Modern Roman}
\usepackage{amsmath,amssymb,mathtools,unicode-math}
\setmathfont{Latin Modern Math}
\usepackage{enumitem,needspace,xurl,hyperref}
\setlength{\parindent}{0pt}
\setlength{\parskip}{5pt}
\setlength{\emergencystretch}{3em}
\newcommand{\N}{\mathbb{N}}
\newcommand{\Z}{\mathbb{Z}}
\newcommand{\Q}{\mathbb{Q}}
\newcommand{\R}{\mathbb{R}}
\newcommand{\C}{\mathbb{C}}
\newcommand{\F}{\mathbb{F}}
\newcommand{\A}{\mathbb{A}}
\newcommand{\PP}{\mathbb{P}}
\newcommand{\E}{\mathbb{E}}
\newcommand{\eps}{\varepsilon}
\newcommand{\dd}{\,\mathrm{d}}
\begin{document}
\section*{UnsolvedMath ID: 30006407}
\textbf{Source problem code:} catalogue entry
\subsection{Definitions and mathematical statement}
Let \(T\) be an area-minimizing current with analytic boundary, as in the supplied source. For a current, the boundary operator is defined by \(\partial T(\omega)=T(\mathrm d\omega)\) on appropriate test forms. Area minimization means minimization of mass among the permitted competitors with the same boundary; the source does not specify the coefficient group, ambient space, current dimension, or whether minimization and comparisons are local or global. The phrase ``analytic boundary'' also requires clarification concerning the analytic geometric support, its multiplicity, and any embeddedness assumptions.

Write \(B_{\partial}(T)\) provisionally for the boundary branching set intended by the source. This notation does not supply a definition: the record gives no local criterion distinguishing a boundary branch point from other boundary singularities. Let \(\mathcal C\) denote the intended class of admissible currents. The extracted research question is to describe the geometric and regularity structure of \(B_{\partial}(T)\) for every \(T\in\mathcal C\), and to determine its sharp dimension bound. Schematically, the latter asks for
\[
\sup_{T\in\mathcal C}\dim_{*} B_{\partial}(T),
\]
with any necessary dependence on current dimension, ambient dimension, codimension, or boundary multiplicity made explicit. Here \(\dim_{*}\) is a placeholder: the source does not specify a notion of dimension, so interpreting it as Hausdorff dimension would be an additional assumption. Neither the class \(\mathcal C\) nor the set \(B_{\partial}(T)\) is sufficiently specified to turn this schematic question into a uniquely quantified mathematical statement.

The supplied context attributes interior branching dimension bounds to Almgren and subsequent work of De Lellis and Spadaro, and contrasts these with the boundary question for analytic boundaries. It supplies neither the precise interior theorem nor a proposed boundary bound. Its dated literature assessment describes the boundary problem as open despite progress in two-dimensional and special boundary settings; these are supplied contextual qualifications, not independently verified mathematical or literature conclusions.
\subsection{Short English statement}
The question asks how branching can occur at the analytic boundary of an area-minimizing current and how large the boundary branching set can be. A precise formulation requires the admissible current class, a definition of boundary branching, the meaning of analytic boundary, and the intended notion of dimension. The supplied open status is dated curation; this is a statement-only draft requiring human mathematical review.
\subsection{Sources}
\begin{itemize}
\item Local UnsolvedMath registry, permanent ID \(30006407\), source problem code \texttt{catalogue entry}, title ``Boundary Branching of Area-Minimizing Currents,'' in the supplied Oberwolfach Reports Open Problems collection. The supplied proposal year is 2025; no individual proposer is identified.
\item Supplied report citation: \emph{Partial Differential Equations (2026)}, Oberwolfach Reports, DOI 10.4171/owr/2025/33, pp. 1753--1822. \url{https://doi.org/10.4171/owr/2025/33}. The literature-source list labels this same reference ``Oberwolfach Report 33/2025.'' Both year descriptions are retained pending verification.
\item Supplied literature reference: \emph{An essential one-sided boundary singularity for a 3-dimensional area-minimizing current}. \url{https://arxiv.org/abs/2502.15047}. No authors or precise theorem from this reference are supplied.
\item The original context credits Almgren and subsequent work of De Lellis and Spadaro for interior branching dimension bounds. No specific bibliographic references for those works are supplied.
\item Literature assessment checked at \texttt{2026-08-22}: the supplied curation labels the problem open and reports progress in two-dimensional and special boundary settings. This is supplied dated curation derived from the OWR material, not an independently verified assessment; time-sensitive status and cited claims require independent verification.
\item The registry labels the statement \texttt{corrected\_verified} and reports that the official report was checked and notation repaired. These are supplied provenance claims; no source visit or independent verification is asserted here.
\end{itemize}
\end{document}
